\section*{Capítulo \ref{ch:g1:shapes} --- Figuras y espacio}

\begin{solution}{exo:g1:shapes:1}
$3$ esquinas: un triángulo. Redonda y sin esquinas: un círculo. $4$ lados
iguales: un cuadrado.
\end{solution}

\begin{solution}{exo:g1:shapes:2}
Las respuestas varían: una puerta y un libro son rectángulos; un plato o un reloj
son círculos; una señal de peligro es un triángulo; una baldosa puede ser un
cuadrado.
\end{solution}

\begin{solution}{exo:g1:shapes:3}
Falso: un cuadrado girado sobre una esquina sigue teniendo cuatro lados iguales y
cuatro esquinas --- girarlo no cambia la figura.
\end{solution}

\begin{solution}{exo:g1:shapes:4}
Dibujos en la cuadrícula: el cuadrado mide $3 \times 3$ casillas, el rectángulo
$5 \times 2$, y el triángulo usa como lado una línea inclinada de la
cuadrícula.
\end{solution}

\begin{solution}{exo:g1:shapes:5}
Las respuestas varían. Por ejemplo: la pelota está a la izquierda, el coche está entre
la pelota y la muñeca, y la muñeca está a la derecha.
\end{solution}

\begin{solution}{exo:g1:shapes:6}
La cruz se dibuja dentro del círculo, la estrella fuera y el punto
justo sobre la línea del círculo.
\end{solution}

\begin{solution}{exo:g1:shapes:7}
Siguiendo $\to\ \uparrow\ \to\ \uparrow\ \to$ avanzas $3$ casillas
a la derecha y $2$ hacia arriba. La vuelta es $\downarrow\ \downarrow$ y luego
$\leftarrow\ \leftarrow\ \leftarrow$ (o cualquier camino que regrese a la
salida).
\end{solution}

\begin{solution}{exo:g1:shapes:8}
Dos triángulos y un cuadrado: $3 + 3 + 4 = 10$ esquinas. Un círculo
y dos rectángulos: $0 + 4 + 4 = 8$ esquinas.
\end{solution}

\begin{solution}{exo:g1:shapes:9}
Con $12$ cerillas sale un cuadrado con $3$ cerillas en cada lado
($3 + 3 + 3 + 3 = 12$). Con $10$ cerillas no se puede: $10$
no se reparte en $4$ partes enteras iguales (cada lado necesitaría
$2$ cerillas y media).
\end{solution}
